% 原书 p.29–p.33（PDF p.37–p.41）
% >>> p037
\section{泊松括号（Poisson Bracket）}

\begin{keytakeaways}
\begin{itemize}
  \item 物理图像：泊松括号给出相空间上两个可观测量沿彼此哈密顿流的相互变化率，
        即 $\{F,G\}=\dd F(X_G)$；它把经典力学的「运动积分」语言代数化。
  \item 数学工具：哈密顿向量场、Lie 导数与内乘的 Cartan 公式、
        斜对称性与 Jacobi 恒等式、$\omega$ 的闭性与非退化性。
  \item 几何观点：$\{F,G\}=\omega(X_G,X_F)$；$\{F,G\}$ 的哈密顿向量场恰为
        $[X_F,X_G]$，故哈密顿向量场在 Lie 括号下封闭；$L$ 是拉格朗日子流形当且仅当
        $I_L=\{F:F|_L=0\}$ 对泊松括号封闭。
  \item 本节结论：引理 2.3.1 给出 $n\,\dd G\wedge\dd F\wedge\omega^{n-1}=\{F,G\}\omega^n$，
        由此得辛结构形变的辛性上界 $1/|\{F,G\}|_{C^0}$；Entov--Polterovich 的
        $C^0$-刚性 (2.14) 及其对多重括号的失效。
\end{itemize}
\end{keytakeaways}

泊松括号是经典力学与哈密顿动力学中的基本概念。它直接处理辛流形上的函数及其相互关系。事实上，在辛几何中，泊松括号常以相当出人意料的方式出现。作为本节的主体，我们将展示若干基于泊松括号的意想不到的性质。

\begin{wdefn}{2.3.1}
设 $(M,\omega)$ 为辛流形。给定两个函数 $F,G\in C^\infty(M)$，定义 $F$ 与 $G$ 的
\textbf{泊松括号}（Poisson bracket）如下
\begin{equation}
  \{F,G\}:=\dd F(X_G),
  \tag{2.8}
\end{equation}
其中 $X_G$ 是 (1.15) 中定义的 $G$ 的哈密顿向量场。
\end{wdefn}
% [Tutor Note: 易错预警——$\{F,G\}$ 的下标顺序：由 $\dd F=-\iota_{X_F}\omega$ 得
% $\{F,G\}=\dd F(X_G)=-\omega(X_F,X_G)=\omega(X_G,X_F)$，即括号中「第二个函数」的哈密顿向量场排在
% $\omega$ 的第一个槽位。若换成 $\sum_i\dd q_i\wedge\dd p_i$ 的辛形式约定，符号与槽位次序都会随之改变，
% 核对时务必与定义 1.7.1 的符号约定联立检查。]

下面是由定义 2.3.1 直接得到的几点有用观察。

(i) 对任意 $F,G\in C^\infty(M)$，有
\begin{equation}
  \{F,G\}=-\dd F(-X_G)=\omega(X_F,-X_G)=\omega(X_G,X_F).
  \tag{2.9}
\end{equation}
由这一性质，$\{F,G\}=-\omega(X_F,X_G)=-\{G,F\}$，即泊松括号是斜对称的。

(ii) 若 $F$ 或 $G$ 为常数，则 $\{F,G\}=0$，因为 $\dd F=0$。事实上，由 (2.8)，只要
$F$ 沿哈密顿向量场 $X_G$ 为常数，就有 $\{F,G\}=0$。于是，由上面的 (i) 可得到一个
不那么显然的结论：$G$ 也沿哈密顿向量场 $X_F$ 为常数（因为
$\{G,F\}=-\{F,G\}=0$）。

(iii) 对任意 $F\in C^\infty(M)$，有 $\{F,F\}=\omega(X_F,X_F)=0$。这既可由上面的
(i) 看出，也可由上面的 (ii) 借助「$F$ 总沿其自身的哈密顿向量场 $X_F$ 为常数」这一
有用事实看出。

(iv) 由 $\{G,F\}$ 生成的哈密顿向量场，即 $X_{\{G,F\}}$，满足如下关系
\begin{equation}
  X_{\{G,F\}}=[X_F,X_G].
  \tag{2.10}
\end{equation}
这直接说明（向量场的）Lie 括号在哈密顿向量场集合中是封闭的。验证 (2.10) 并非平凡，
故我们在下面给出细节。

% >>> p038
回顾：对 $M$ 上任意 2-形式 $\omega$，其沿向量场 $a$ 的 Lie 导数定义为
\begin{equation}
  (\mathcal{L}_a\omega)(b,c)
  =a\cdot\omega(b,c)-\omega([a,b],c)-\omega(b,[a,c]),
  \tag{2.11}
\end{equation}
其中 $b,c$ 是 $M$ 上的向量场。这里第一项 $a\cdot\omega(b,c)$ 表示函数 $\omega(b,c)$
沿向量场 $a$ 的方向导数。类似地，对 $M$ 上任意 1-形式 $\lambda$，有
\begin{equation}
  (\mathcal{L}_a\lambda)(b)=a\cdot\lambda(b)-\lambda([a,b]).
  \tag{2.12}
\end{equation}
利用 (2.11) 与 (2.12)，可算得
\begin{align*}
  (\mathcal{L}_a\iota_b\omega-\iota_b\mathcal{L}_a\omega)(c)
  &=a\cdot\omega(b,c)-\omega(b,[a,c])-a\cdot\omega(b,c)+\omega([a,b],c)
    +\omega(b,[a,c])\\
  &=\omega([a,b],c)=(\iota_{[a,b]}\omega)(c).
\end{align*}
现在取 $a=X_F$、$b=X_G$，得
\begin{align*}
  \iota_{[X_F,X_G]}\omega
  &=\mathcal{L}_{X_F}\iota_{X_G}\omega-\iota_{X_G}\mathcal{L}_{X_F}\omega\\
  &=\dd(\iota_{X_F}\iota_{X_G}\omega)+\iota_{X_G}\dd(\iota_{X_F}\omega)\\
  &=\dd(\omega(X_G,X_F))=-\dd\{G,F\}.
\end{align*}
这里，对第二个等式与第三个等式，我们用到 $\omega$ 闭这一条件，加上 $\iota_{X_G}\omega$
与 $\iota_{X_F}\omega$ 都是恰当的（从而特别地是闭的）；对最后一个等式，我们应用 (2.9)。
因此，由定义 1.7.1，$-\dd\{G,F\}=\iota_{X_{\{G,F\}}}\omega$，便得到 (2.10) 中欲证的结论。
% [Tutor Note: 上述第二个等式来自 Cartan 公式 $\mathcal{L}_X=\dd\iota_X+\iota_X\dd$ 作用于
% $\iota_{X_G}\omega$：$\mathcal{L}_{X_F}\iota_{X_G}\omega=\dd\iota_{X_F}\iota_{X_G}\omega
% +\iota_{X_F}\dd(\iota_{X_G}\omega)$；而 $\iota_{X_G}\omega=-\dd G$ 恰当，故
% $\dd(\iota_{X_G}\omega)=0$，该项整体为零。同理 $\mathcal{L}_{X_F}\omega=0$ 使
% $\iota_{X_G}\mathcal{L}_{X_F}\omega=0$。原文此处的两项写法只差一个不影响结果的零项，
% 核对时抓住「恰当形式的 $\dd$ 为零」即可。]

注意 (2.8) 中泊松括号的定义只依赖于 $\omega$ 的非退化性（定义 1.1.1 中的条件 (i)），
因为 (2.8) 与 (1.15) 中哈密顿向量场 $X_G$ 的存在性正来自 $\omega$ 的非退化性。下面这个
练习说明，从 $\omega$ 的闭性（定义 1.1.1 中的条件 (ii)）能推出什么。

\begin{wexer}{2.3.1}
（[171] 中的定理 2.2.9）证明：非退化 2-形式 $\omega$ 是闭的，当且仅当泊松括号满足
如下 \emph{Jacobi 恒等式}
\[
  \{F,\{G,H\}\}+\{G,\{H,F\}\}+\{H,\{F,G\}\}=0
\]
对任意光滑函数 $F,G,H$ 成立。
\end{wexer}

接下来我们给出辛几何中两个泊松括号意外出现的例子。

\begin{wexam}{2.3.1}
（辛结构的形变，见 [35] 第 1.6 节）给定紧辛流形 $(M,\omega)$，可对 $\omega$ 作形变，
使形变后的 2-形式仍是辛的。讨论任意形变会很困难，故为简单起见，我们只考虑如下情形。
固定 $F,G\in C^\infty(M)$，考虑一个恰当形变（即附加项是恰当形式）：
\[
  \omega_s:=\omega-s\cdot\dd G\wedge\dd F\qquad\text{对某个 } s\ge0 .
\]
% >>> p039
显然 $\omega_s$ 是闭的；此外，由于 $\omega$ 的非退化性是一个开条件，当 $s$ 充分小时，
形变后的 $\omega_s$ 也非退化（从而 $\omega_s$ 是辛的）。这里有一个自然的问题。\hfill$//$
\end{wexam}

\begin{wque}{2.3.1}
使 $\omega_s$ 保持辛性的 $s$ 的上界是什么？
\end{wque}

为回答这个问题，我们需要如下引理。

\begin{wlem}{2.3.1}
设 $(M,\omega)$ 是 $2n$ 维辛流形。则对任意 $F,G\in C^\infty(M)$，有
\[
  n\cdot\dd G\wedge\dd F\wedge\omega^{n-1}=\{F,G\}\cdot\omega^n .
\]
\end{wlem}

\begin{wproof}
考虑形式 $\dd F\wedge\omega^n$（它恒等于零，因为它是 $M^{2n}$ 上的 $(2n+1)$-形式），
于是有
\begin{align*}
  0&=\iota_{X_G}(\dd F\wedge\omega^n)
  =\{F,G\}\cdot\omega^n-\dd F\wedge(\iota_{X_G}\omega^n)\\
  &=\{F,G\}\cdot\omega^n-n\cdot\dd F\wedge(\iota_{X_G}\omega)\wedge\omega^{n-1}\\
  &=\{F,G\}\cdot\omega^n+n\cdot\dd F\wedge\dd G\wedge\omega^{n-1}.
\end{align*}
于是得到欲证的等式。
\end{wproof}

于是由引理 2.3.1，有
\[
  (\omega_s)^n=\omega^n-s\cdot n\cdot\dd G\wedge\dd F\wedge\omega^{n-1}
  =(1-s\{F,G\})\omega^n .
\]
因此当 $s\in[0,1/(|\{F,G\}|_{C^0}))$ 时，$\omega_s$ 仍是辛的。这里 $|\cdot|_{C^0}$ 表示
$M$ 上的「上确界范数」。换言之，$1/(|\{F,G\}|_{C^0})$ 给出了问题 2.3.1 中所要求的
上界。
% [Tutor Note: 这里用到了两点：(1) 二项式展开 $(\omega-s\,\dd G\wedge\dd F)^n$ 中凡含
% $(\dd G\wedge\dd F)^2$ 的项全部消失（因为 $(\dd G\wedge\dd F)^2=0$），只剩
% $\omega^n$ 与 $n$ 个 $\dd G\wedge\dd F$ 的交叉项；(2) $\omega_s$ 非退化等价于
% $(\omega_s)^n$ 处处非零，故只需 $1-s\{F,G\}\neq0$ 且 $s\ge0$。
% 注意上界用 $\{F,G\}$ 的 $C^0$ 范数而非最大值，是因为 $\{F,G\}$ 可能取负值。]

下面这个例子揭示泊松括号与拉格朗日子流形之间的联系。更确切地说，它从函数的角度提供了
判定一个子流形是否为拉格朗日子的判别准则。当判定 $L$ 的局部切空间有困难时，这一准则
尤为有用。

\begin{wexam}{2.3.2}
（拉格朗日子的泊松判别法）设 $L\subset(M^{2n},\omega)$ 表示一个 $n$ 维子流形。考虑
$C^\infty(M)$ 的如下子空间
\[
  I_L:=\{F\in C^\infty(M)\mid F|_L=0\}.
\]
则 $L$ 是拉格朗日子流形，当且仅当泊松括号在 $I_L$ 中封闭，即对任意 $F,G\in I_L$ 都有
$\{F,G\}|_L=0$。\footnote{第一作者 2017 年在耶路撒冷希伯来大学 D. Kazhdan 的
「Sunday seminar」上首次得知这一结果。}

事实上，设 $L$ 是拉格朗日子流形。对任意 $p\in L$，
\begin{equation}
  \{F,G\}(p)=\dd F(X_G(p))
  =\left.\frac{\dd}{\dd t}F\big(\varphi_G^t(p)\big)\right|_{t=0},
  \tag{2.13}
\end{equation}
% >>> p040
其中 $\varphi_G^t$ 是由 $X_G$ 生成的哈密顿流。由于 $G|_L=0$，可知 $X_G\in TL$（见
命题 1.6.1 的证明）。故 $\varphi_G^t(p)\in L$。然而 $F\in I_L$，故 $F(\varphi_G^t(p))=0$，
由 (2.13) 即得 $\{F,G\}(p)=0$。

反之，设 $L$ 不是拉格朗日的，即存在点 $p\in L$ 与向量 $v\in T_pM$，使得
$\omega_p(v,-)\equiv0$ 在 $T_pL$ 上成立，但 $v\notin T_pL$。取函数 $G\in C^\infty(M)$
使得
\[
  G|_L=0\qquad\text{但}\qquad \dd G_p(v)\neq0 .
\]
这样的 $G$ 可局部构造（例如借助定理 1.2.1 中的 Darboux 定理）。于是对由 $G$ 生成的
哈密顿向量场 $X_G$，有
\[
  \omega_p(X_G(p),v)=-\omega_p(v,X_G(p))=\dd G_p(v)\neq0,
\]
% [Tutor Note: 原文疑误——按定义 1.7.1 的约定 $\dd G=-\iota_{X_G}\omega$，故
% $\dd G_p(v)=-\omega_p(X_G(p),v)$，上式末端的等号应带负号，即
% $\omega_p(X_G(p),v)=-\dd G_p(v)$。因该式的用途只是说明「$\neq0$」，符号不影响论证。
% 译文保留原文写法，详见 review_ch02.md。]
所以 $X_G(p)$ 不在 $T_pL$ 中。用同样的办法可构造 $F\in C^\infty(M)$，使得 $F|_L=0$
而 $\dd F_p(X_G(p))\neq0$。于是由定义 2.3.1，有
$\{F,G\}(p)=\dd F_p(X_G(p))\neq0$。因此 $\{F,G\}$ 不属于 $I_L$，这与「泊松括号在
$I_L$ 中封闭」的假设矛盾。\hfill$//$
\end{wexam}

为结束本节，我们提及辛几何中的 $C^0$-刚性现象。这里 $C^0$-刚性是指：对象在
$C^0$-拓扑下收敛到一个受限制的极限。最著名的一个是如下事实（归功于 Gromov$^{[100]}$
与 Y. Eliashberg$^{[73]}$）：$\Symp(M,\omega)$ 在 $\Diff(M)$ 中是 $C^0$-闭的。更确切
地说，对 $(M,\omega)$ 上一列辛同胚 $\{\varphi_k\}_{k\in\N}$，若它 $C^0$-收敛到一个
微分同胚 $\varphi$，则极限 $\varphi$ 也必定是辛同胚！这之所以令人惊讶，是因为按定义，
要验证 $\varphi$ 是辛同胚，至少需要考虑它的导数（而这已经属于 $C^1$ 范畴的信息）。
然而一般说来，$C^0$-收敛\emph{不能}控制 $C^1$ 信息。

对泊松括号，我们确实有类似的 $C^0$-刚性，但较弱。

\begin{wthm}{2.3.1}
（Entov--Polterovich [80]）对辛流形 $(M,\omega)$ 上的任意一对函数
$F,G\in C_c^\infty(M)$，有
\begin{equation}
  |\{F,G\}|_{C^0}
  =\liminf_{F',G'\xrightarrow{C^0}F,G}|\{F',G'\}|_{C^0}.
  \tag{2.14}
\end{equation}
\end{wthm}
% [Tutor Note: (2.14) 中的 $\liminf$ 取遍所有在 $C^0$-拓扑下分别收敛到 $F,G$ 的函数对
% $(F',G')$（不必有界于 $C^1$）。等式的含义是：泊松括号的 $C^0$ 范数不能靠 $C^0$-扰动
% 突然变大——它恰是「$C^0$ 下半连续」的精确量化。注意 $C_c^\infty$ 表示紧支光滑函数，
% 紧支性用于控制 Hofer 范数（见证明第 2 步）。]

作为该定理证明的提纲，我们列出如下几个步骤。

第一步：设 $H,K\in C_c^\infty(M)$ 是不含时的哈密顿函数，分别生成哈密顿同痕
$\{\varphi_H^t\}_{t\in[0,1]}$ 与 $\{\varphi_K^t\}_{t\in[0,1]}$。考虑换位子
$\{\varphi_H^t\varphi_K^t\varphi_H^{-t}\varphi_K^{-t}\}_{t\in[0,1]}$，它由函数
$L_t=H-H\circ\varphi_K^{-1}\circ\varphi_H^{-t}$ 生成。特别地，取 $H=rF$、$K=sG$，
可把 $L_t$ 重写为
\[
  L_{r,s}^t:=rF-rF\circ(\varphi_G^s)^{-1}\circ(\varphi_F^{rt})^{-1}.
\]
% [Tutor Note: 「换位子生成函数」的公式来自 Hamilton 流的复合公式：若 $\varphi^t$ 由 $H$ 生成，
% 则 $\psi^{-1}\circ\varphi^t\circ\psi$ 由 $H\circ\psi$ 生成；而两个流复合生成的函数需用
% 流形上的「Magri--Morosi」型公式 $H\#K=H+K\circ\varphi_H^{-1}$ 递推。此处 $L_t$ 的符号须
% 与定义 1.7.1 的约定一致，核对时建议先用 $H,K$ 为线性函数、流为平移的玩具情形验证。
% 另注：原文此处记为 $L_{r,s}^t$，与下一行 $\varphi_{L_{r,s}^1}$ 的下标约定须对应阅读。]
% >>> p041
第二步：在第 6 章我们将看到，哈密顿微分同胚群 $\Ham(M,\omega)$ 上存在一个伪范数
$\|\cdot\|_{\mathrm{Hofer}}$。类似地，可定义该范数的「正部」，记作
$\|\cdot\|_{\mathrm{Hofer}}^+$，如下
\[
  \|\varphi\|_{\mathrm{Hofer}}^+
  :=\inf_{H_t}\int_0^1\max_MH_t\,\dd t,
\]
其中下确界取遍所有生成连接 $\mathbf{1}$ 与 $\varphi\in\Ham(M,\omega)$ 的路径
$\{\varphi_t\}_{t\in[0,1]}$ 的哈密顿函数 $H_t$。于是，借助
$\|\cdot\|_{\mathrm{Hofer}}^+$（作为泛函）的连续性，可得：对任意小 $\epsilon>0$，
存在充分小的 $r,s>0$ 使得
\begin{equation}
  \begin{aligned}
    rs\big(\max\{F,G\}-2\epsilon\big)
    &\le\big\|\varphi_{L_{r,s}^1}\big\|_{\mathrm{Hofer}}^+\\
    &\le rs\Big(\liminf_{F',G'\xrightarrow{C^0}F,G}\max\{F',G'\}+\epsilon\Big).
  \end{aligned}
  \tag{2.15}
\end{equation}
在 (2.15) 两端除以 $rs$ 并令 $\epsilon\to0$，得
\begin{equation}
  \max\{F,G\}\le\liminf_{F',G'\xrightarrow{C^0}F,G}\max\{F',G'\}.
  \tag{2.16}
\end{equation}
第三步：不等式 (2.16) 的反方向易于验证，故有
\begin{equation}
  \max\{F,G\}=\liminf_{F',G'\xrightarrow{C^0}F,G}\max\{F',G'\}.
  \tag{2.17}
\end{equation}
在 (2.17) 中以 $-F$ 替换 $F$，可得到关于 $-\min\{F,G\}$ 的类似结论，由此即得欲证的
等式 (2.14)。
% [Tutor Note: (2.15) 中 $\max\{F,G\}$ 指函数 $\{F,G\}$ 在 $M$ 上的最大值（不是 $F$ 与 $G$
% 逐点取大），这一点极易误读；同理 $\min\{F,G\}$ 是泊松括号的最小值。第 3 步之所以要换成
% $-F$：因为 $\max(-f)=-\min f$，于是对 $-\{F,G\}$ 的估计给出 $\{F,G\}$ 下界一侧的不等式，
% 两端合并才得到 $|\{F,G\}|_{C^0}$ 的等式 (2.14)。]

然而，这一刚性定理对多重泊松括号失效。

\begin{wthm}{2.3.2}
（Entov--Polterovich [80]）对任意辛流形 $(M,\omega)$，都可找到函数
$F,G,H\in C_c^\infty(M)$ 满足 $\{F,\{G,H\}\}\not\equiv0$；但存在光滑函数
$F',G',H'\in C_c^\infty(M)$，分别在 $C^0$ 意义下任意接近 $F,G,H$，并满足条件
\[
  \{F',\{G',H'\}\}\equiv0 .
\]
\end{wthm}

此外，下面这个练习说明泊松括号通常的 $C^0$-连续性在一般情况下失效。

\begin{wexer}{2.3.2}
记 $|F|_{C^0}$ 为函数 $F$ 的上确界范数。找出函数 $F,G\in C^\infty(\R^2)$ 以及函数列
$F_k\xrightarrow{C^0}F$、$G_k\xrightarrow{C^0}G$，使得
\[
  \limsup_{F',G'\xrightarrow{C^0}F,G}|\{F_k,G_k\}|_{C^0}>|\{F,G\}|_{C^0}.
\]
% [Tutor Note: 原文此处下标写作「$F',G'\xrightarrow{C^0}F,G$」而括号内为 $\{F_k,G_k\}$，
% 记号自相矛盾：既然正文给出的是函数列 $F_k\to F$、$G_k\to G$，$\limsup$ 应取在 $k\to\infty$
% 上（或写成 $\limsup_k$）。译文保留原文写法，详见 review_ch02.md。此练习的实质是：泊松括号
% 只对 $C^1$ 收敛连续，$C^0$ 收敛不足以控制导数，故 $\limsup$ 可严格大于极限值。]
\end{wexer}
